\chapter{Evaluation}
\label{ch:evaluation}

\section{Experimental Setup}

Experiments were run on a five-node cluster. \autoref{tab:results}
summarises throughput under increasing partition rates.

\begin{table}[htbp]
  \centering
  \caption{Throughput (requests/s) under simulated partition rates.}
  \label{tab:results}
  \begin{tabular}{@{}lrr@{}}
    \toprule
    Partition rate & Baseline & Proposed \\
    \midrule
    0\%  & 12{,}400 & 12{,}100 \\
    5\%  & 9{,}800  & 11{,}600 \\
    10\% & 6{,}200  & 10{,}900 \\
    \bottomrule
  \end{tabular}
\end{table}

\section{Latency Under Load}

\autoref{fig:latency} shows median request latency as the partition
rate increases.

\begin{figure}[htbp]
  \centering
  \rule{0.6\textwidth}{4cm}
  \caption{Median latency versus simulated partition rate.}
  \label{fig:latency}
\end{figure}

\section{Discussion}

The proposed protocol maintains throughput more effectively than the
baseline once the partition rate exceeds 5\%, at the cost of a small
constant latency overhead \citep{ongaro2014search}.
